% GENERATED FILE. Edit canonical lessons, then run npm run generate:books.
% canonical-source-hash: fnv1a64:1d22c9404579c6b4
% canonical-lessons: FA-C182-rezerv-kardan, FA-C182-estekhdam-kardan, FA-C182-tekrar-kardan, FA-C182-tarjome-kardan, FA-C182-sohbat-kardan

\chapter{To book, To hire, To repeat, To translate, To talk}
\label{ch:fa-to-book-to-hire-to-repeat-to-translate-to-talk}
\hlchaptermodality{182}

\begin{chapteropening}
\textbf{By the end of this chapter:} \emph{I can say to book, to hire, to repeat, to translate and to talk.}

Complete the last lesson of chapter 182: I can say to book, to hire, to repeat, to translate and to talk.
\end{chapteropening}

\section[\texorpdfstring{\fa{رزرو} \fa{کردن} (rezerv kardan)}{رزرو کردن (rezerv kardan)}]{\fa{رزرو} \fa{کردن} (rezerv kardan) — to book, to reserve}

\label{lesson:FA-C182-rezerv-kardan}



\begin{quote}
Before the new one: say the Persian for to cross, then the Persian for to stroll.
\end{quote}

\subsection*{You'll want to know: \fa{رزرو} \fa{کردن}}
\textbf{\fa{رزرو} \fa{کردن}} — \emph{rezerv kardan} — \textquotedblleft{}to book, to reserve\textquotedblright{}.

\begin{cousinweb}
\emph{rezerv kardan} pairs \emph{rezerv}, borrowed from French \emph{réserve}, \textquotedblleft{}reserve\textquotedblright{}, with the Persian verb \emph{kardan}, \textquotedblleft{}to do\textquotedblright{}. \emph{kardan}, \textquotedblleft{}to do\textquotedblright{}, is an inherited verb, from Middle Persian \emph{kardan} (EtymDB).
\end{cousinweb}

\begin{grammarlens}[title={What you've built}]
One more word for this chapter.
\end{grammarlens}

\subsection*{Your turn}
\begin{itemize}
\raggedright
  \item \emph{Say it:} \emph{rezerv kardan}
  \item \emph{Say it:} \emph{rezerv kardan}, once more
  \item \emph{Say it:} say \emph{rezerv kardan}
  \item \emph{Recall:} say the Persian for to cross, then the Persian for to stroll, then say \emph{rezerv kardan} again
\end{itemize}

\begin{tcolorbox}[breakable,colback=teal!4,colframe=teal!35!black,title={Before you move on}]
What does \fa{رزرو} \fa{کردن} mean? (\textquotedblleft{}To book, to reserve\textquotedblright{}.) Say it once more.
\end{tcolorbox}
\section[\texorpdfstring{\fa{استخدام} \fa{کردن} (estekhdâm kardan)}{استخدام کردن (estekhdâm kardan)}]{\fa{استخدام} \fa{کردن} (estekhdâm kardan) — to hire, to employ}

\label{lesson:FA-C182-estekhdam-kardan}



\begin{quote}
Before the new one: say the Persian for to stroll, then the Persian for to book.
\end{quote}

\subsection*{You'll want to know: \fa{استخدام} \fa{کردن}}
\textbf{\fa{استخدام} \fa{کردن}} — \emph{estekhdâm kardan} — \textquotedblleft{}to hire, to employ\textquotedblright{}.

\begin{cousinweb}
\emph{estekhdâm kardan} pairs \emph{estekhdâm}, borrowed from Arabic \emph{istiḫdām}, \textquotedblleft{}taking into service\textquotedblright{} (root \emph{ḫ-d-m}, Lane), with the Persian verb \emph{kardan}, \textquotedblleft{}to do, to make\textquotedblright{}. \emph{kardan}, \textquotedblleft{}to do\textquotedblright{}, is an inherited verb, from Middle Persian \emph{kardan} (EtymDB).
\end{cousinweb}

\begin{grammarlens}[title={What you've built}]
One more word for this chapter.
\end{grammarlens}

\subsection*{Your turn}
\begin{itemize}
\raggedright
  \item \emph{Say it:} \emph{estekhdâm kardan}
  \item \emph{Say it:} \emph{estekhdâm kardan}, once more
  \item \emph{Say it:} say \emph{estekhdâm kardan}
  \item \emph{Recall:} say the Persian for to stroll, then the Persian for to book, then say \emph{estekhdâm kardan} again
\end{itemize}

\begin{tcolorbox}[breakable,colback=teal!4,colframe=teal!35!black,title={Before you move on}]
What does \fa{استخدام} \fa{کردن} mean? (\textquotedblleft{}To hire, to employ\textquotedblright{}.) Say it once more.
\end{tcolorbox}
\section[\texorpdfstring{\fa{تکرار} \fa{کردن} (tekrâr kardan)}{تکرار کردن (tekrâr kardan)}]{\fa{تکرار} \fa{کردن} (tekrâr kardan) — to repeat}

\label{lesson:FA-C182-tekrar-kardan}



\begin{quote}
Before the new one: say the Persian for to book, then the Persian for to hire.
\end{quote}

\subsection*{You'll want to know: \fa{تکرار} \fa{کردن}}
\textbf{\fa{تکرار} \fa{کردن}} — \emph{tekrâr kardan} — \textquotedblleft{}to repeat\textquotedblright{}.

\begin{cousinweb}
\emph{tekrâr kardan} pairs \emph{tekrâr}, borrowed from Arabic \emph{takrār}, \textquotedblleft{}repetition\textquotedblright{} (root \emph{k-r-r}, Lane), with the Persian verb \emph{kardan}, \textquotedblleft{}to do, to make\textquotedblright{}. \emph{kardan}, \textquotedblleft{}to do\textquotedblright{}, is an inherited verb, from Middle Persian \emph{kardan} (EtymDB).
\end{cousinweb}

\begin{grammarlens}[title={What you've built}]
One more word for this chapter.
\end{grammarlens}

\subsection*{Your turn}
\begin{itemize}
\raggedright
  \item \emph{Say it:} \emph{tekrâr kardan}
  \item \emph{Say it:} \emph{tekrâr kardan}, once more
  \item \emph{Say it:} say \emph{tekrâr kardan}
  \item \emph{Recall:} say the Persian for to book, then the Persian for to hire, then say \emph{tekrâr kardan} again
\end{itemize}

\begin{tcolorbox}[breakable,colback=teal!4,colframe=teal!35!black,title={Before you move on}]
What does \fa{تکرار} \fa{کردن} mean? (\textquotedblleft{}To repeat\textquotedblright{}.) Say it once more.
\end{tcolorbox}
\section[\texorpdfstring{\fa{ترجمه} \fa{کردن} (tarjome kardan)}{ترجمه کردن (tarjome kardan)}]{\fa{ترجمه} \fa{کردن} (tarjome kardan) — to translate}

\label{lesson:FA-C182-tarjome-kardan}



\begin{quote}
Before the new one: say the Persian for to hire, then the Persian for to repeat.
\end{quote}

\subsection*{You'll want to know: \fa{ترجمه} \fa{کردن}}
\textbf{\fa{ترجمه} \fa{کردن}} — \emph{tarjome kardan} — \textquotedblleft{}to translate\textquotedblright{}.

\begin{cousinweb}
\emph{tarjome kardan} pairs \emph{tarjome}, borrowed from Arabic \emph{tarjama}, \textquotedblleft{}translation\textquotedblright{} (root \emph{t-r-j-m}, Lane), with the Persian verb \emph{kardan}, \textquotedblleft{}to do, to make\textquotedblright{}. \emph{kardan}, \textquotedblleft{}to do\textquotedblright{}, is an inherited verb, from Middle Persian \emph{kardan} (EtymDB). Arabic itself took the word from Aramaic (EtymDB).
\end{cousinweb}

\begin{grammarlens}[title={What you've built}]
One more word for this chapter.
\end{grammarlens}

\subsection*{Your turn}
\begin{itemize}
\raggedright
  \item \emph{Say it:} \emph{tarjome kardan}
  \item \emph{Say it:} \emph{tarjome kardan}, once more
  \item \emph{Say it:} say \emph{tarjome kardan}
  \item \emph{Recall:} say the Persian for to hire, then the Persian for to repeat, then say \emph{tarjome kardan} again
\end{itemize}

\begin{tcolorbox}[breakable,colback=teal!4,colframe=teal!35!black,title={Before you move on}]
What does \fa{ترجمه} \fa{کردن} mean? (\textquotedblleft{}To translate\textquotedblright{}.) Say it once more.
\end{tcolorbox}
\section[\texorpdfstring{\fa{صحبت} \fa{کردن} (sohbat kardan)}{صحبت کردن (sohbat kardan)}]{\fa{صحبت} \fa{کردن} (sohbat kardan) — to talk, to chat}

\label{lesson:FA-C182-sohbat-kardan}



\begin{quote}
Before the new one: say the Persian for to repeat, then the Persian for to translate.
\end{quote}

\subsection*{You'll want to know: \fa{صحبت} \fa{کردن}}
\textbf{\fa{صحبت} \fa{کردن}} — \emph{sohbat kardan} — \textquotedblleft{}to talk, to chat\textquotedblright{}.

\begin{cousinweb}
\emph{sohbat kardan} pairs \emph{sohbat}, borrowed from Arabic \emph{ṣuḥba}, \textquotedblleft{}companionship\textquotedblright{} (root \emph{ṣ-ḥ-b}, Lane), with the Persian verb \emph{kardan}, \textquotedblleft{}to do, to make\textquotedblright{}. \emph{kardan}, \textquotedblleft{}to do\textquotedblright{}, is an inherited verb, from Middle Persian \emph{kardan} (EtymDB). Persian moved it from keeping company to talking (EtymDB).
\end{cousinweb}

\begin{grammarlens}[title={What you've built}]
One more word for this chapter.
\end{grammarlens}

\subsection*{Your turn}
\begin{itemize}
\raggedright
  \item \emph{Say it:} \emph{sohbat kardan}
  \item \emph{Say it:} \emph{sohbat kardan}, once more
  \item \emph{Say it:} say \emph{sohbat kardan}
  \item \emph{Recall:} say the Persian for to repeat, then the Persian for to translate, then say \emph{sohbat kardan} again
\end{itemize}

\begin{tcolorbox}[breakable,colback=teal!4,colframe=teal!35!black,title={Before you move on}]
What does \fa{صحبت} \fa{کردن} mean? (\textquotedblleft{}To talk, to chat\textquotedblright{}.) Say it once more.
\end{tcolorbox}
